% =====================================================================
% band-theme.tex
% Identidad visual Centro BAND para Quarto + beamer (tema metropolis).
% Paleta y logos tomados de "Plantilla BAND presentaciones.pptx".
%
% Uso desde un .qmd:
%   format:
%     beamer:
%       theme: metropolis
%       include-in-header:
%         - file: band-theme.tex
%
% Requiere la carpeta images/ (logos y fondos) junto al .qmd.
% =====================================================================

\usepackage{booktabs}
\usepackage{multirow}
\usepackage{amsmath}
\usepackage{bm}
\usepackage{etoolbox}
\usepackage{tikz}
\usetikzlibrary{arrows.meta,positioning}
\graphicspath{{images/}}

% ---------------------------------------------------------------------
% 1. Paleta BAND (colores del tema del .pptx)
% ---------------------------------------------------------------------
\definecolor{band1}{HTML}{007A8C}      % petróleo (color principal)
\definecolor{band2}{HTML}{239C9F}      % verde azulado
\definecolor{band3}{HTML}{00C2BB}      % turquesa
\definecolor{band4}{HTML}{00CC99}      % verde
\definecolor{bandred}{HTML}{E6605B}    % coral
\definecolor{bandorange}{HTML}{F26B43} % naranja
\definecolor{bandyellow}{HTML}{D7F49A} % lima
\definecolor{bandlila}{HTML}{8F94F7}   % lila
\definecolor{bandblue}{HTML}{3700EF}   % azul eléctrico
\definecolor{bandpink}{HTML}{F8D3F3}   % rosado
\definecolor{bandcream}{HTML}{EEECE3}  % crema (fondo del .pptx)
\definecolor{bandink}{HTML}{0E2841}    % azul noche (texto)

% ---------------------------------------------------------------------
% 2. Colores y fuentes de beamer / metropolis
% ---------------------------------------------------------------------
\setbeamercolor{normal text}{fg=bandink,bg=white}
\setbeamercolor{background canvas}{bg=white}
\setbeamercolor{structure}{fg=band1}

% Título de lámina: texto petróleo sobre franja crema.
% (Para la franja sólida del template anterior: fg=white,bg=band1)
\setbeamercolor{frametitle}{fg=band1,bg=bandcream}
\setbeamerfont{frametitle}{size=\Large,series=\bfseries}

% Barra de progreso de metropolis
\setbeamercolor{progress bar}{fg=band3,bg=band3!25}
\makeatletter
\setlength{\metropolis@progressinheadfoot@linewidth}{1.2pt}
\makeatother

% Énfasis y listas
\setbeamercolor{alerted text}{fg=bandred}
\setbeamercolor{example text}{fg=band2}
\setbeamercolor{itemize item}{fg=band2}
\setbeamercolor{itemize subitem}{fg=band3}
\setbeamercolor{enumerate item}{fg=band1}
\setbeamercolor{enumerate subitem}{fg=band2}
\setbeamerfont{enumerate item}{series=\bfseries}
\setbeamercovered{transparent=20}
\setbeamertemplate{frametitle continuation}{}
\hypersetup{colorlinks=true,allcolors=band1}

% Bloques (encabezados #### dentro de una lámina)
%   #### Título            -> petróleo
%   #### Título {.example} -> verde azulado
%   #### Título {.alert}   -> coral
\setbeamercolor{block title}{bg=band1,fg=white}
\setbeamercolor{block body}{bg=band1!8,fg=bandink}
\setbeamercolor{block title example}{bg=band2,fg=white}
\setbeamercolor{block body example}{bg=band2!10,fg=bandink}
\setbeamercolor{block title alerted}{bg=bandred,fg=white}
\setbeamercolor{block body alerted}{bg=bandred!10,fg=bandink}
\setbeamerfont{block title}{series=\bfseries,size=\normalsize}
\setbeamerfont{block title example}{series=\bfseries,size=\normalsize}
\setbeamerfont{block title alerted}{series=\bfseries,size=\normalsize}

% Código un poco más pequeño para que quepa en columnas
\AtBeginEnvironment{Shaded}{\small}
\AtBeginEnvironment{verbatim}{\small}

% ---------------------------------------------------------------------
% 3. Callouts de Quarto con la paleta BAND
%    note      -> petróleo  (conceptos, definiciones)
%    tip       -> verde     (buenas prácticas)
%    warning   -> naranja   (errores frecuentes)
%    important -> coral     (reglas que no se rompen)
%    caution   -> azul      (ejercicios; usar title="Ejercicio")
% ---------------------------------------------------------------------
\makeatletter
\AtBeginDocument{%
  \definecolor{quarto-callout-note-color}{HTML}{007A8C}%
  \definecolor{quarto-callout-note-color-frame}{HTML}{007A8C}%
  \definecolor{quarto-callout-tip-color}{HTML}{00CC99}%
  \definecolor{quarto-callout-tip-color-frame}{HTML}{00CC99}%
  \definecolor{quarto-callout-warning-color}{HTML}{F26B43}%
  \definecolor{quarto-callout-warning-color-frame}{HTML}{F26B43}%
  \definecolor{quarto-callout-important-color}{HTML}{E6605B}%
  \definecolor{quarto-callout-important-color-frame}{HTML}{E6605B}%
  \definecolor{quarto-callout-caution-color}{HTML}{3700EF}%
  \definecolor{quarto-callout-caution-color-frame}{HTML}{3700EF}%
  % Texto de los callouts alineado a la izquierda (sin justificar)
  \@ifpackageloaded{tcolorbox}{\tcbset{halign=flush left}}{}%
}
\makeatother

% ---------------------------------------------------------------------
% 4. Pie de página: isotipo + título corto | sección | n / N + logo
%    Cambiar el título corto desde el .qmd con:
%      \renewcommand{\bandshorttitle}{Mi título corto}
% ---------------------------------------------------------------------
\newcommand{\bandshorttitle}{\insertshorttitle}
\newlength{\bandfootmargin}
\setlength{\bandfootmargin}{0.2cm}   % espacio a cada lado del pie de página

\setbeamertemplate{footline}{%
	\ifnum\value{page}>1\relax
	\begin{beamercolorbox}[wd=\paperwidth,ht=0.52cm,dp=0.22cm]{footline}%
		\hspace*{\bandfootmargin}%
		\makebox[\dimexpr0.37\paperwidth-\bandfootmargin\relax][l]{%
			%        \raisebox{-0.13cm}{\includegraphics[height=0.42cm]{band_isotipo.png}}\hspace{0.2cm}%
			{\scriptsize\textcolor{band1}{\bandshorttitle}}}%
		\makebox[0.26\paperwidth][c]{\scriptsize\textcolor{band2}{\insertsection}}%
		\makebox[\dimexpr0.37\paperwidth-\bandfootmargin\relax][r]{%
			{\scriptsize\textcolor{band1}{\insertframenumber\,/\,\inserttotalframenumber}}%
			\hspace{0.3cm}%
			\raisebox{-0.09cm}{\includegraphics[height=0.36cm]{band_logo.png}}}%
	\end{beamercolorbox}%
	\fi
}
% ---------------------------------------------------------------------
% 5. Fondos especiales: portada, portadilla de sección, destacado, cierre
%    El fondo se elige con una marca que cada lámina activa desde su
%    contenido y que se borra al comenzar la lámina siguiente.
% ---------------------------------------------------------------------
\makeatletter
\newcommand{\band@bg}{}
\newcommand{\band@bgcolor}{band1}
\newcommand{\band@closingtext}{}
\AtBeginEnvironment{frame}{\gdef\band@bg{}}

% Lienzo del tamaño de la lámina; coordenadas de (0,0) a (1,1)
\newcommand{\band@canvas}[1]{%
  \begin{tikzpicture}[x=\paperwidth,y=\paperheight,inner sep=0pt,outer sep=0pt]
    \useasboundingbox (0,0) rectangle (1,1);
    #1
  \end{tikzpicture}}

% --- Portada (siempre la página 1) ---
\newcommand{\band@drawportada}{\band@canvas{%
  \node[anchor=south west] at (0,0)
    {\includegraphics[width=\paperwidth,height=\paperheight]{band_fondo_portada.png}};
  \node[anchor=north west] at (0.035,0.945)
    {\includegraphics[height=1.5cm]{band_logo_bajada_blanco.png}};
  \node[anchor=north east,text=white,font=\scriptsize] at (0.965,0.925) {\insertdate};
  \node[anchor=west,text=white,text width=0.40\paperwidth,align=flush left] at (0.56,0.5) {%
    {\LARGE\bfseries\inserttitle\par}%
    \ifx\insertsubtitle\@empty\else\vspace{0.25cm}{\large\insertsubtitle\par}\fi
    \vspace{0.3cm}{\color{bandyellow}\rule{1.2cm}{1.5pt}\par}\vspace{0.3cm}%
    {\small\insertauthor\par}%
    \ifx\insertinstitute\@empty\else\vspace{0.08cm}{\scriptsize\insertinstitute\par}\fi};
  \node[anchor=south west] at (0.035,0.06)
    {\includegraphics[height=0.75cm]{band_isotipo_claro.png}};
  \node[anchor=south east] at (0.965,0.075)
    {\includegraphics[height=0.3cm]{band_web_blanco.png}};
}}

% --- Portadilla de sección ---
\newcommand{\band@drawseccion}{\band@canvas{%
  \node[anchor=south west] at (0,0)
    {\includegraphics[width=\paperwidth,height=\paperheight]{band_fondo_crema.png}};
  \node[anchor=north west] at (0.035,0.945)
    {\includegraphics[height=0.8cm]{band_logo.png}};
  \node[anchor=west,text width=0.40\paperwidth,align=flush left] at (0.56,0.5) {%
    {\fontsize{30}{32}\selectfont\bfseries\color{band3}%
      \ifnum\value{section}<10 0\fi\arabic{section}\par}%
    \vspace{0.2cm}{\color{bandred}\rule{1.2cm}{1.5pt}\par}\vspace{0.3cm}%
    {\LARGE\bfseries\color{band1}\insertsection\par}};
  \node[anchor=south west] at (0.035,0.06)
    {\includegraphics[height=0.75cm]{band_isotipo.png}};
  \node[anchor=south east] at (0.965,0.075)
    {\includegraphics[height=0.3cm]{band_web.png}};
}}

% --- Portadilla de subsección ---
% Color de fondo: cambiar aquí (tonos claros: bandpink, bandyellow, bandcream)
%\newcommand{\bandsubseccioncolor}{band4}
%\newcommand{\band@drawsubseccion}{\band@canvas{%
%		\fill[color=\bandsubseccioncolor] (0,0) rectangle (1,1);
%		\node[anchor=north west] at (0.035,0.945)
%		{\includegraphics[height=0.8cm]{band_logo.png}};
%		\node[anchor=west,text width=0.40\paperwidth,align=flush left] at (0.56,0.5) {%
%			{\small\bfseries\color{band1}%
%				\ifnum\value{section}<10 0\fi\arabic{section}\hspace{0.6em}\insertsection\par}%
%			\vspace{0.2cm}{\color{bandred}\rule{1.2cm}{1.5pt}\par}\vspace{0.3cm}%
%			{\Large\bfseries\color{band1}\insertsubsection\par}};
%		\node[anchor=south west] at (0.035,0.06)
%		{\includegraphics[height=0.75cm]{band_isotipo.png}};
%		\node[anchor=south east] at (0.965,0.075)
%		{\includegraphics[height=0.3cm]{band_web.png}};
%}}


% --- Portadilla de subsección (fondo verde con la "B") ---
\newcommand{\band@drawsubseccion}{\band@canvas{%
		\node[anchor=south west] at (0,0)
		{\includegraphics[width=\paperwidth,height=\paperheight]{band_fondo_verde.png}};
		\node[anchor=north west] at (0.035,0.945)
		{\includegraphics[height=0.8cm]{band_logo_petroleo.png}};
		  \node[anchor=west,text width=0.40\paperwidth,align=flush left] at (0.56,0.5) {%
			{\small\bfseries\color{band1}%
				\ifnum\value{section}<10 0\fi\arabic{section}\hspace{0.6em}\insertsection\par}%
			\vspace{0.2cm}{\color{bandred}\rule{1.2cm}{1.5pt}\par}\vspace{0.3cm}%
			{\Large\bfseries\color{band1}\insertsubsection\par}};
		\node[anchor=south west] at (0.035,0.06)
		{\includegraphics[height=0.75cm]{band_isotipo.png}};
		\node[anchor=south east] at (0.965,0.075)
		{\includegraphics[height=0.3cm]{band_web.png}};
}}



% --- Lámina destacada a color completo ---
\newcommand{\band@drawcolor}{\band@canvas{%
  \fill[color=\band@bgcolor] (0,0) rectangle (1,1);
  \node[anchor=south west] at (0.035,0.06)
    {\includegraphics[height=0.75cm]{band_isotipo_claro.png}};
  \node[anchor=south east] at (0.965,0.075)
    {\includegraphics[height=0.3cm]{band_web_blanco.png}};
}}

% --- Cierre ---
\newcommand{\band@drawcierre}{\band@canvas{%
  \fill[color=bandcream] (0,0) rectangle (1,1);
  \node at (0.5,0.60) {\includegraphics[height=2.6cm]{band_logo_bajada.png}};
  \node at (0.5,0.34) {\includegraphics[height=0.42cm]{band_web.png}};
  \node[text=band1,font=\small,align=center] at (0.5,0.20) {\band@closingtext};
}}

\setbeamertemplate{background}{%
  \ifnum\value{page}=1\relax
    \band@drawportada
  \else
    \ifdefstring{\band@bg}{seccion}{\band@drawseccion}{%
	\ifdefstring{\band@bg}{subseccion}{\band@drawsubseccion}{%
		\ifdefstring{\band@bg}{color}{\band@drawcolor}{%
			\ifdefstring{\band@bg}{cierre}{\band@drawcierre}{}}}}%
  \fi
}

% La portada se dibuja completa en el fondo (título, autor, fecha, logos)
\setbeamertemplate{title page}{\mbox{}}

% Comandos para usar dentro de una lámina {.plain .noframenumbering}
%   \bandseccion                    portadilla de sección (automática con #)
%   \banddestacado[color]{Texto}    mensaje clave; color: band1 (defecto),
%                                   band2, bandred, bandblue, bandlila
%   \bandcierre[Texto de contacto]  lámina final con logo
\newcommand{\bandseccion}{\gdef\band@bg{seccion}\mbox{}}
\newcommand{\bandsubseccion}{\gdef\band@bg{subseccion}\mbox{}}
\newcommand{\banddestacado}[2][band1]{%
  \gdef\band@bg{color}\gdef\band@bgcolor{#1}%
  \begin{center}\color{white}\LARGE\bfseries #2\par\end{center}}
\newcommand{\bandcierre}[1][]{%
  \gdef\band@bg{cierre}\gdef\band@closingtext{#1}\mbox{}}
\makeatother

% Portadilla automática al inicio de cada sección (#)
\AtBeginSection[]{%
  \begin{frame}[plain,noframenumbering]
    \bandseccion
  \end{frame}}
\AtBeginSubsection[]{%
	\begin{frame}[plain,noframenumbering]
		\bandsubseccion
\end{frame}}

% ---------------------------------------------------------------------
% 6. Marcador de posición para capturas de pantalla
%    \bandcaptura{Texto}  o  \bandcaptura[0.7]{Texto}  (alto/ancho)
% ---------------------------------------------------------------------
\newcommand{\bandcaptura}[2][0.6]{%
  \begin{center}
  \begin{tikzpicture}
    \node[draw=band2,dashed,line width=0.6pt,rounded corners=3pt,fill=bandcream,
          minimum width=0.96\linewidth,minimum height=#1\linewidth,
          text width=0.8\linewidth,align=center,text=band1,font=\small,
          inner sep=0pt] {#2};
  \end{tikzpicture}
  \end{center}}
